% from ega2-1-fr.tex
\oldpage[II]{34}

est commutatif. Il suffit encore de le vérifier dans chaque $D_+(f)$ et
cela découle aussitôt des définitions et de
(\hyperref[I.1.3.13-fr]{I, 1.3.13}).
\end{env}

\begin{env}[2.5.12]
\label{II.2.5.12-fr}
Sous les hypothèses de (\hyperref[II.2.5.11-fr]{2.5.11}), on définit un
homomorphisme canonique fonctoriel de $S_{(f)}$-modules
\[
\label{II.2.5.12.1-fr}
  \mu_f:\bigl(\operatorname{Hom}_S(M,N)\bigr)_{(f)}
  \longrightarrow
  \operatorname{Hom}_{S_{(f)}}(M_{(f)},N_{(f)})
\tag{2.5.12.1}
\]
en faisant correspondre à $u/f^n$, où $u$ est un homomorphisme de degré
$nd$, l'homomorphisme $M_{(f)}\to N_{(f)}$ qui transforme $x/f^m$
($x\in M_{md}$) en $u(x)/f^{m+n}$. Pour $g\in S_e$ ($e>0$), on a encore
un diagramme commutatif
\[
  \xymatrix{
    \bigl(\operatorname{Hom}_S(M,N)\bigr)_{(f)}
    \ar[r]^{\mu_f}\ar[d]
    & \operatorname{Hom}_{S_{(f)}}(M_{(f)},N_{(f)})\ar[d]\\
    \bigl(\operatorname{Hom}_S(M,N)\bigr)_{(fg)}
    \ar[r]_{\mu_{fg}}
    & \operatorname{Hom}_{S_{(fg)}}(M_{(fg)},N_{(fg)})
  }
\]
(la flèche de gauche étant l'homomorphisme canonique, celle de droite
provenant des homomorphismes canoniques). On en conclut encore (compte
tenu de (\hyperref[I.1.3.8-fr]{I, 1.3.8})) que les $\mu_f$ définissent un
homomorphisme canonique fonctoriel de $\mathcal{O}_X$-Modules
\[
\label{II.2.5.12.2-fr}
  \mu:\widetilde{\operatorname{Hom}_S(M,N)}
  \longrightarrow
  \mathcal{H}om_{\mathcal{O}_X}(\widetilde{M},\widetilde{N}).
\tag{2.5.12.2}
\]
\end{env}

\begin{proposition}[2.5.13]
\label{II.2.5.13-fr}
Supposons l'idéal $S_+$ engendré par $S_1$. Alors l'homomorphisme
$\lambda$ (\hyperref[II.2.5.11.2-fr]{2.5.11.2}) est un isomorphisme~; il
en est de même de l'homomorphisme $\mu$
(\hyperref[II.2.5.12.2-fr]{2.5.12.2}) lorsque le $S$-module gradué $M$
admet une présentation finie (\hyperref[II.2.1.1-fr]{2.1.1}).
\end{proposition}

\begin{proof}
Comme $X$ est réunion des $D_+(f)$ pour $f\in S_1$
(\hyperref[II.2.3.14-fr]{2.3.14}), on est ramené à prouver que $\lambda_f$
et $\mu_f$ sont des isomorphismes, sous les hypothèses envisagées, lorsque
$f$ est homogène et \emph{de degré $1$}. Or, on définit alors une
application $\mathbf{Z}$-bilinéaire
\[
  M_m\times N_n\longrightarrow
  M_{(f)}\otimes_{S_{(f)}}N_{(f)}
\]
en faisant correspondre à $(x,y)$ l'élément
$(x/f^m)\otimes(y/f^n)$ (lorsque $m<0$, nous écrivons $x/f^m$ pour
$f^{-m}x/1$)~; ces applications définissent une application
$\mathbf{Z}$-linéaire
\[
  M\otimes_{\mathbf{Z}}N\longrightarrow
  M_{(f)}\otimes_{S_{(f)}}N_{(f)},
\]
et si $s\in S_q$, cette application transforme $(sx)\otimes y$ en
\[
  (s/f^q)\bigl((x/f^m)\otimes(y/f^n)\bigr)
\]
(pour $x\in M_m$, $y\in N_n$). On en déduit donc un di-homomorphisme de
modules
\[
  \gamma_f:M\otimes_S N\longrightarrow
  M_{(f)}\otimes_{S_{(f)}}N_{(f)},
\]
relatif à l'homomorphisme canonique $S\to S_{(f)}$ (transformant
$s\in S_q$ en $s/f^q$). Supposons en outre que pour un élément
$\sum_i(x_i\otimes y_i)$ de $M\otimes_S N$ ($x_i$, $y_i$ homogènes de
degrés respectifs $m_i$, $n_i$) on ait
\[
  f^r\sum_i(x_i\otimes y_i)=0,
\]
autrement dit $\sum_i(f^rx_i\otimes y_i)=0$. On en déduit en vertu de
(\hyperref[0.1.3.4-fr]{0, 1.3.4}) que l'on a
\[
  \sum_i(f^rx_i/f^{m_i+r})\otimes(y_i/f^{n_i})=0,
\]
c'est-à-dire $\gamma_f(\sum_i(x_i\otimes y_i))=0$. Par suite $\gamma_f$
se factorise en
\[
  M\otimes_S N\longrightarrow(M\otimes_S N)_f
  \xrightarrow{\gamma'_f}
  M_{(f)}\otimes_{S_{(f)}}N_{(f)}~;
\]
si $\lambda'_f$ est la restriction de $\gamma'_f$
